% Scope: Connect nonuniform sampling to gridding and operator-based reconstruction; explain acceleration as a joint sampling and prior-design problem.
% Status: architecture only; no chapter prose or completed problems yet.
% Prerequisites: Chapters 5, 6 and 14; solver details in Chapter 16.
% Planned worked examples: Radial sampling density, spiral off-resonance and adjoint checks.
% Planned figures: Trajectories, interpolation kernels and undersampling PSFs.
% Planned challenge problems: Design a trajectory under slew constraints; distinguish density compensation from inversion.
\chapter{Non-Cartesian and Accelerated Imaging}
\label{ch:15}

\section{Radial and spiral acquisition}
\label{sec:15:radial-and-spiral-acquisition}
\subsection{Angular sampling and golden-angle ordering}
\subsection{Spiral geometry and variable density}
\subsection{Stack-of-stars and three-dimensional trajectories}

\section{Nonuniform Fourier reconstruction}
\label{sec:15:nonuniform-fourier-reconstruction}
\subsection{Gridding kernels and oversampling}
\subsection{Deapodization and density compensation}
\subsection{NUFFT forward and adjoint operators}

\section{Trajectory and field errors}
\label{sec:15:trajectory-and-field-errors}
\subsection{Gradient delays and measured trajectories}
\subsection{Off-resonance blur and time segmentation}
\subsection{Motion robustness and self-navigation limits}

\section{Sampling for acceleration}
\label{sec:15:sampling-for-acceleration}
\subsection{Variable-density masks and incoherence}
\subsection{Partial Fourier with complex phase}
\subsection{Joint coil and sampling design}

\section{Dynamic imaging}
\label{sec:15:dynamic-imaging}
\subsection{Temporal interleaving and view sharing}
\subsection{Spatiotemporal undersampling and low-rank structure}
\subsection{Temporal footprint versus nominal frame rate}

\section{Worked examples}
\label{sec:15:worked-examples}
% Planned calculations: Radial sampling density, spiral off-resonance and adjoint checks.
\subsection{Derivation and quantitative calculation}
\subsection{Assumptions, limiting cases and scanner interpretation}

\section{Challenge problems}
\label{sec:15:challenge-problems}
% Problem design: Design a trajectory under slew constraints; distinguish density compensation from inversion.
% Use challengeproblem and stable labels prob:15:<topic>.
% Complete solutions belong in Appendix E, section for this chapter.
% Use \begin{solution}{prob:15:<topic>} ... \end{solution} there.
\subsection{Derivation and quantitative design}
\subsection{Model criticism and integrated reasoning}
